\chapter{Метод серверно-сліпого каналу нагляду}
\label{ch:relay}

% ============================================================
%  РОЗДІЛ 7 -- privacy-preserving data-plane метод.
%  Джерело: aura-platform (canonical) + aura-messenger-ios.
%  ПОТРЕБА: зафіксований відтворюваний E2E-доказ (див. розділ 8/додатки).
% ============================================================

\section{Постановка задачі: приватний контрольований перегляд}
\label{sec:relay-problem}
% -- Дитину треба захищати, але сервер не має бачити приватний safety payload.

\section{Згода, контрольована child endpoint-ом}
\label{sec:relay-consent}
% -- child-controlled consent як умова активації каналу.

\section{Активний parent-child зв'язок}
\label{sec:relay-binding}
% -- parent-child binding; авторизаційна модель.

\section{Provisioning pair-key для кожного пристрою}
\label{sec:relay-pairkey}
% -- Початковий relay pair key генерується на клієнті,
%    запаковується per parent device через наявний E2E key-package шлях.

\section{Маршрутизація opaque payload і роль сервера}
\label{sec:relay-routing}
% -- Сервер виконує лише authz, rate limiting, зберігання ciphertext,
%    доставку sealed key packages; pair key і payload не читає.

\section{Життєвий цикл ACK / revoke / purge}
\label{sec:relay-lifecycle}
% -- ACK лише після parent install; revoke/purge -> crypto-shredding стану.

\section{Межі: replay, TTL, DoS та межа приватності}
\label{sec:relay-limits}
% -- Контрзаходи; privacy boundary.

\section{Наскрізний сценарій та відтворюваний доказ}
\label{sec:relay-e2e}
% -- grant consent -> list recipient devices -> wrap pair key -> submit package ->
%    parent consume -> ACK -> child sends sealed relay message -> parent opens -> revoke purges.
Метод ілюструється наскрізним сценарієм від надання згоди до відкликання, що супроводжується відтворюваним інтеграційним тестом і артефактним документом з точною командою та очікуваним результатом (див.\ розділ~8 і додатки). \emph{[Зафіксувати E2E-доказ.]}

\rozdilconcl{7}
\emph{[Висновки до розділу 7.]}
